\section{{\fontsize{12pt}{14pt}\selectfont RESULTADOS}}
\vspace{-0.3cm}

    
\novoBlocoTexto{
    Os resultados são apresentados em duas partes. A seção~\ref{sec:secao-5-1} apresenta a plataforma desenvolvida e o atendimento aos requisitos levantados na seção~\ref{sec:levantamento-requisitos}, e as seções~\ref{sec:secao-5-2} a~\ref{sec:discussao-limitacoes} apresentam a execução e os resultados da avaliação experimental descrita na seção~\ref{sec:avaliacao-experimental}.
}

\subsection{{\fontsize{12pt}{14pt}\selectfont A plataforma desenvolvida}}
\label{sec:secao-5-1}
\novoBlocoTexto{
    A plataforma foi implementada como uma aplicação web, acessível por meio de navegadores (RNF08). Para representar o processo de utilização da plataforma, foi empregada a notação \textit{Business Process Model and Notation} (BPMN), utilizada para modelar e representar graficamente processos. O fluxo foi organizado em três diagramas BPMN: preparação da disciplina (Figura~\ref{fig:bpmn-preparacao}), envio e processamento das dúvidas (Figura~\ref{fig:bpmn-duvidas}) e consulta ao painel docente (Figura~\ref{fig:bpmn-painel}). Os diagramas apresentam as principais atividades e interações entre professor, aluno e sistema.
}

\bpmnPreparacao

\bpmnDuvidas

\bpmnPainel

\novoBlocoTexto{
    As Figuras de~\ref{fig:tela-acesso} a~\ref{fig:tela-topico-trechos} apresentam suas telas na ordem do fluxo descrito no capítulo~\ref{sec:metodo}: o acesso à plataforma, a disciplina preparada pelo professor, o envio de uma dúvida pelo aluno e a consulta do professor ao painel. O Quadro~\ref{quad:atendimento-requisitos}, ao final desta seção, relaciona cada requisito dos Quadros~\ref{quad:requisitos-funcionais} e~\ref{quad:requisitos-nao-funcionais} à figura que comprova o seu atendimento.
}

\novoBlocoTexto{
    O cadastro e a autenticação são os mesmos para professores e alunos (Figura~\ref{fig:tela-acesso}). 
}
\telaacesso

\novoBlocoTexto{
    O papel de cada usuário é definido na disciplina: o professor a cadastra, e a plataforma gera os códigos de convite com os quais alunos e outros professores ingressam nela (Figura~\ref{fig:tela-disciplina}). 

\teladisciplina
    Ao acessar a página de tópicos e materiais da disciplina (Figura~\ref{fig:tela-topico}) o professor envia o plano de ensino em PDF, a LLM identifica os tópicos e gera uma breve descrição de cada um, que o professor pode editar antes de salvar. Em seguida, o professor associa a cada tópico os materiais em PDF, dos quais a plataforma extrai o texto, divide-o em fragmentos e gera os \textit{embeddings} de cada fragmento. A descrição extraída do plano aparece no início da página de cada tópico, como em Pilhas (Figura~\ref{fig:tela-topico-resumos}), e os arquivos associados aparecem como origem dos slides recuperados (Figura~\ref{fig:tela-topico-trechos}).
}

\telatopico

\novoBlocoTexto{
    O aluno envia suas dúvidas pelo \textit{chat} e recebe a resposta em português (RNF06), elaborada pela LLM com base nos fragmentos recuperados do material e acompanhada da fonte e do slide de onde veio cada informação, como “[Fonte 2, slide 11]” (Figura~\ref{fig:tela-chat}). A classificação e o resumo da dúvida são gerados no mesmo processamento e não aparecem para o aluno. Abaixo da caixa de mensagem, o \textit{chat} informa ao aluno que o professor poderá visualizar resumos das dúvidas enviadas, sem vínculo à identificação do estudante (RNF09), como mostra a Figura~\ref{fig:tela-chat}.
}

\telachat

\novoBlocoTexto{
    Ao acessar o painel, o professor vê o número de disciplinas, o total de dúvidas classificadas, o número de tópicos com dúvidas, a lista dos tópicos com mais dúvidas e um gráfico com a distribuição das dúvidas entre os tópicos. A Figura~\ref{fig:tela-painel} mostra esse painel ao final do experimento, com as 1.416 perguntas classificadas nos 12 tópicos da disciplina.
}

\telapainel

\novoBlocoTexto{
    Ao abrir um tópico, o professor vê a sua descrição, o número de perguntas registradas e, para cada dúvida, um título e um resumo individual, escritos em terceira pessoa, a pergunta original e sem a identificação do aluno (RNF01), como mostra a Figura~\ref{fig:tela-topico-resumos} para o tópico Pilhas. Na mesma página, uma segunda aba apresenta os slides mais recuperados pela busca nas dúvidas do tópico, com o arquivo de origem de cada um (Figura~\ref{fig:tela-topico-trechos}), o que indica ao professor as partes do material mais recuperadas para a resolução das dúvidas.
}

\telatopicoresumos
\telatopicotrechos

\novoBlocoTexto{
    O controle de acesso é aplicado no próprio banco de dados, por políticas de segurança em nível de linha (RLS) que restringem cada conversa e as suas mensagens ao aluno que as criou (RNF02 e RNF03). Além disso, cada rota da aplicação confere se o usuário participa da disciplina, e as instruções enviadas à LLM tratam a dúvida e os fragmentos como dados, com a orientação de ignorar comandos contidos neles.
}

\novoBlocoTexto{
    O Quadro \ref{quad:atendimento-requisitos} resume o atendimento aos requisitos. Todos foram atendidos. O aviso do RNF09 foi incluído na tela do \textit{chat} depois do experimento, o que não altera os resultados, pois as perguntas foram enviadas por automação (seção~\ref{sec:execucao-classificacao}).
}

\begin{quadro}[H]\captionsetup{name=Quadro}
\centering
\caption{Atendimento aos Requisitos}
\label{quad:atendimento-requisitos}

\renewcommand{\arraystretch}{1.3}
\setlength{\tabcolsep}{4pt}

\begin{tabularx}{\textwidth}{
|>{\raggedright\arraybackslash}X
|>{\raggedright\arraybackslash}X|
}
\hline

\rowcolor[HTML]{C0C0C0}
\multicolumn{1}{|c|}{\textbf{Requisito}} &
\multicolumn{1}{c|}{\textbf{Evidência}} \\ \hline

RF01 -- Cadastro de disciplinas &
Figuras~\ref{fig:tela-disciplina} e~\ref{fig:tela-painel} \\ \hline

RF02 -- Convite de alunos por código de acesso &
Figura~\ref{fig:tela-disciplina} \\ \hline

RF03 -- Envio do plano de ensino e extração dos tópicos &
Figura~\ref{fig:tela-topico-resumos} e Seção~\ref{sec:secao-5-2} \\ \hline

RF04 -- Materiais em PDF associados aos tópicos &
Figura~\ref{fig:tela-topico-trechos} e Seção~\ref{sec:secao-5-2} \\ \hline

RF05 -- Extração do texto dos materiais &
Seção~\ref{sec:secao-5-2} \\ \hline

RF06 -- Divisão do texto em fragmentos &
Seção~\ref{sec:secao-5-2} \\ \hline

RF07 -- Geração dos \textit{embeddings} de cada fragmento &
Seção~\ref{sec:secao-5-2} \\ \hline

RF08 -- Envio de dúvidas pelo \textit{chat} &
Figura~\ref{fig:tela-chat} \\ \hline

RF09 -- Recuperação dos cinco fragmentos &
Figuras~\ref{fig:tela-chat} e~\ref{fig:tela-topico-trechos} \\ \hline

RF10 -- Resposta com base nos fragmentos &
Figura~\ref{fig:tela-chat} \\ \hline

RF11 -- Classificação em um dos tópicos &
Figura~\ref{fig:tela-painel} e Seção~\ref{sec:secao-5-3} \\ \hline

RF12 -- Resumo individual de cada dúvida &
Figura~\ref{fig:tela-topico-resumos} \\ \hline

RF13 -- Armazenamento da pergunta, do tópico e do resumo &
Seção~\ref{sec:secao-5-2} \\ \hline

RF14 -- Painel com quantidades e resumos por tópico &
Figuras~\ref{fig:tela-painel} e~\ref{fig:tela-topico-resumos} \\ \hline

RNF01 -- Professor sem acesso aos nomes &
Figura~\ref{fig:tela-topico-resumos} \\ \hline

RNF02 -- Aluno com acesso só às próprias conversas &
Seção~\ref{sec:secao-5-1} \\ \hline

RNF03 -- Controle de acesso por RLS &
Seção~\ref{sec:secao-5-1} \\ \hline

RNF04 -- LLM executada localmente &
Quadro~\ref{quad:configuracao-ambiente-experimental} \\ \hline

RNF05 -- \textit{Embeddings} com suporte ao português &
Quadro~\ref{quad:configuracao-ambiente-experimental} \\ \hline

RNF06 -- Respostas em português &
Figura~\ref{fig:tela-chat} \\ \hline

RNF07 -- Vetores no banco, com \textit{pgvector} &
Quadro~\ref{quad:configuracao-ambiente-experimental} \\ \hline

RNF08 -- Acesso pelo navegador &
Figuras~\ref{fig:tela-acesso},
\ref{fig:tela-disciplina},
\ref{fig:tela-chat},
\ref{fig:tela-painel},
\ref{fig:tela-topico-resumos} e~\ref{fig:tela-topico-trechos} \\ \hline

RNF09 -- Aviso ao aluno sobre o acesso às perguntas &
Figura~\ref{fig:tela-chat}\, \\ \hline

\end{tabularx}
{\fontsize{10pt}{12pt}\selectfont
\caption*{Fonte: Elaborado pela autora (2026).} % Legenda de baixo.
}
\end{quadro}

\FloatBarrier

\subsection{{\fontsize{12pt}{14pt}\selectfont Execução do experimento}}
\label{sec:secao-5-2}
\novoBlocoTexto{
    Na preparação da disciplina, a plataforma extraiu do plano de ensino os 12 tópicos da Tabela~\ref{tab:distribuicao-de-perguntas-por-topico-no-conjunto-de-referencia}, cada um com a sua descrição, e os 12 arquivos de slides da Tabela~\ref{tab:distribuicao-materiais-por-topico}, com 405 slides no total, foram associados aos respectivos tópicos, divididos em fragmentos e convertidos em embeddings (RF03 a RF07).
}

\novoBlocoTexto{
    As 1.416 perguntas traduzidas foram enviadas pelo \textit{chat} a partir de uma única conta de aluno, criada para o teste; por isso, a página de cada tópico registra um único aluno com dúvidas (Figura \ref{fig:tela-topico-resumos}). Todas as perguntas foram classificadas em um dos 12 tópicos da disciplina, e o painel registrou 1.416 dúvidas classificadas, uma para cada pergunta do conjunto (Figura~\ref{fig:tela-painel}). Para cada pergunta, o banco de dados armazenou a pergunta, o tópico atribuído, o grau de confiança, os fragmentos recuperados e o resumo (RF13).
}

\novoBlocoTexto{
    O Quadro \ref{quad:configuracao-ambiente-experimental} apresenta o ambiente de execução. A LLM e o modelo de embeddings foram executados no próprio computador, sem envio das dúvidas a serviços externos de inferência (RNF04), e os vetores foram armazenados no banco de dados da aplicação, com a extensão pgvector (RNF07).
}



\begin{quadro}[H]\captionsetup{name=Quadro}
\centering
\caption{Configuração do ambiente experimental}
\label{quad:configuracao-ambiente-experimental}

\renewcommand{\arraystretch}{1.3}
\setlength{\tabcolsep}{4pt}

\begin{tabularx}{\textwidth}{
|>{\raggedright\arraybackslash}p{3.5cm}
|>{\raggedright\arraybackslash}X|
}
\hline

\rowcolor[HTML]{C0C0C0}
\multicolumn{1}{|c|}{\textbf{Item}} &
\multicolumn{1}{c|}{\textbf{Configuração}} \\ \hline

Computador &
\textit{Notebook} Acer Nitro AN517-54 \\ \hline

Processador &
Intel Core i7-11800H, de 11ª geração, com 8 núcleos e 16 \textit{threads} \\ \hline

Memória RAM &
64 GB \\ \hline

Placa de vídeo &
NVIDIA GeForce RTX 3050 Laptop GPU \\ \hline

Sistema operacional &
Ubuntu 26.04.1 LTS, 64 bits \\ \hline

LLM &
\textit{Qwen2.5-7B-Instruct}, quantizada em \texttt{q4\_k\_m} (formato GGUF), executada pelo \textit{llama.cpp} na placa de vídeo \\ \hline

Temperatura da LLM &
0,2 \\ \hline

Modelo de \textit{embeddings} &
\texttt{intfloat/multilingual-e5-base} (768 dimensões), executado no processador por um serviço local em Python, com \textit{FastAPI} e \textit{sentence-transformers} \\ \hline

Banco de dados &
\textit{Supabase}, em nuvem (\textit{PostgreSQL} com a extensão \textit{pgvector}) \\ \hline

Aplicação &
\textit{Next.js} 16.3.5, \textit{React} 19.2.8 e \textit{TypeScript} 5 \\ \hline

Envio automatizado &
\textit{Selenium WebDriver} 4.52.0, pelo \textit{chat} da plataforma no navegador \\ \hline

\end{tabularx}

{\fontsize{10pt}{12pt}\selectfont
\caption*{Fonte: Elaborado pela autora (2026).} % Legenda de baixo.
}
\end{quadro}

\FloatBarrier

\FloatBarrier

\novoBlocoTexto{
    Ao final da execução, os registros das perguntas foram exportados do banco de dados, com o tópico atribuído pela LLM, e comparados aos tópicos de referência em planilha eletrônica. A planilha com as 1.416 perguntas traduzidas, o tópico de referência e o tópico atribuído pela plataforma está disponível no \textit{Google Drive}\footnote{Disponível em: \url{https://docs.google.com/spreadsheets/d/12FplYIpg0O_vLE_PhnKJ0lyG4CDuXaGG/edit?usp=sharing}.}, serviço de armazenamento de arquivos em nuvem do Google.
}

\subsection{{\fontsize{12pt}{14pt}\selectfont Desempenho da classificação}}
\label{sec:secao-5-3}
\novoBlocoTexto{
    Das 1.416 perguntas do conjunto de referência, a plataforma atribuiu 1.180 ao mesmo tópico da referência, o que corresponde a uma acurácia de 83,33\%. Esse valor é quatro vezes o parâmetro mínimo definido na seção~\ref{sec:metricas-analise-erros}, de 20,8\%, obtido ao atribuir todas as perguntas ao tópico mais frequente. Todas as perguntas foram classificadas em um dos 12 tópicos da disciplina, e a Tabela~\ref{tab:desempenho-classificacao-topico} apresenta o desempenho em cada um deles, as médias macro e a acurácia geral.
}



\begin{table}[htbp]
\centering
\caption{Desempenho da classificação por tópico}
\label{tab:desempenho-classificacao-topico}

\begingroup
\scriptsize
\renewcommand{\arraystretch}{1.3}
\setlength{\tabcolsep}{2pt}

\begin{tabularx}{\textwidth}{
>{\raggedright\arraybackslash}X
>{\centering\arraybackslash}p{2.0cm}
>{\centering\arraybackslash}p{2.5cm}
>{\centering\arraybackslash}p{1.8cm}
>{\centering\arraybackslash}p{1.5cm}
>{\centering\arraybackslash}p{1.8cm}
>{\centering\arraybackslash}p{1.8cm}
}
\hline

\textbf{Tópico} &
\textbf{Perguntas na referência} &
\textbf{Atribuídas pela plataforma} &
\textbf{Acertos} &
\textbf{Precisão (\%)} &
\textbf{Revocação (\%)} &
\textbf{F1-score (\%)} \\ \hline

Revisão de Python & 295 & 137 & 135 & 98,54 & 45,76 & 62,50 \\

Listas & 21 & 85 & 19 & 22,35 & 90,48 & 35,85 \\

Ponteiros & 15 & 36 & 15 & 41,67 & 100,00 & 58,82 \\

Nós & 14 & 8 & 5 & 62,50 & 35,71 & 45,45 \\

Lista simplesmente encadeada & 32 & 47 & 29 & 61,70 & 90,63 & 73,42 \\

Lista duplamente encadeada & 105 & 121 & 102 & 84,30 & 97,14 & 90,27 \\

Pilhas & 196 & 169 & 163 & 96,45 & 83,16 & 89,32 \\

Filas & 24 & 33 & 21 & 63,64 & 87,50 & 73,68 \\

Recursão & 163 & 206 & 160 & 77,67 & 98,16 & 86,72 \\

\textit{Hashing} (Espalhamento) & 163 & 164 & 156 & 95,12 & 95,71 & 95,41 \\

Árvore Binária de Busca & 249 & 262 & 242 & 92,37 & 97,19 & 94,72 \\

Grafo & 139 & 148 & 133 & 89,86 & 95,68 & 92,68 \\ \hline

\textbf{Média macro} &
&
&
&
\textbf{73,85} &
\textbf{84,76} &
\textbf{74,90} \\

\textbf{Total} &
\textbf{1.416} &
\textbf{1.416} &
\textbf{1.180} &
&
&
\\

\textbf{Acurácia geral (\%)} &
\multicolumn{6}{c}{\textbf{83,33} (1.180 acertos em 1.416 perguntas)} \\ \hline

\end{tabularx}
\endgroup
{\fontsize{10pt}{12pt}\selectfont
\caption*{Fonte: Elaborado pela autora (2026).}
}

\end{table}

\FloatBarrier

\novoBlocoTexto{
    O desempenho variou entre os tópicos. \textit{Hashing} (Espalhamento), Árvore binária de busca, Grafo, Lista duplamente encadeada e Pilhas obtiveram F1-score acima de 89\%, enquanto Listas, Nós, Ponteiros e Revisão de Python ficaram entre 35,85\% e 62,50\%. O macro F1-score, calculado pela média simples dos F1-scores de cada tópico e atribuindo o mesmo peso a todos eles (seção~\ref{sec:metricas-analise-erros}), foi de 74,90\%, abaixo da acurácia geral de 83,33\%.
}
\novoBlocoTexto{
    Em dois grupos de tópicos, precisão e revocação seguiram direções opostas. Revisão de Python teve a maior precisão (98,54\%) e uma das menores revocações (45,76\%): quase todas as perguntas atribuídas a esse tópico eram dele, mas 160 das suas 295 perguntas foram atribuídas a outros. Listas e Ponteiros mostraram o inverso, com revocação de 90,48\% e 100\% e precisão de 22,35\% e 41,67\%, porque receberam perguntas de outros tópicos, sobretudo de Revisão de Python. A Figura~\ref{fig:matriz-confusao} apresenta a matriz de confusão, com o destino de cada pergunta.
}

\matrizconfusao

\FloatBarrier

\subsection{{\fontsize{12pt}{14pt}\selectfont Análise dos erros}}
\label{sec:analise-erros}
\novoBlocoTexto{
    Das 236 perguntas atribuídas a um tópico diferente do de referência, 160 (67,8\%) eram de Revisão de Python e 33 (14,0\%) de Pilhas; os outros dez tópicos somaram 43 (18,2\%). Sem as perguntas de Revisão de Python, a concordância nos outros onze tópicos foi de 93,22\% (1.045 de 1.121). Os pares mais frequentes na matriz de confusão foram Revisão de Python → Listas (63 perguntas), Revisão de Python → Recursão (35), Revisão de Python → Ponteiros (11), Revisão de Python → Lista simplesmente encadeada (10) e Pilhas → Recursão (10).
}

\novoBlocoTexto{
    A leitura das perguntas do conjunto aponta três causas prováveis. A primeira é a divisão diferente dos tópicos. No conjunto, Revisão de Python reúne perguntas sobre recursos da linguagem, e 58 das suas 295 perguntas citam listas, vetores (\textit{arrays}), matrizes ou tuplas. No material da disciplina, as listas do Python são o conteúdo do tópico Listas, para o qual a plataforma levou 63 perguntas de Revisão de Python, 39 delas entre essas 58. As outras 24 tratam de outros recursos da linguagem, como laços de repetição, tipos mutáveis e imutáveis, cadeias de caracteres (\textit{strings}) e dicionários, e parte desses assuntos também aparece no material de Listas, que apresenta a lista como sequência mutável e percorre os seus elementos com a função \texttt{enumerate}.
}

\novoBlocoTexto{
    A segunda é que parte dos rótulos parece indicar o tema da aula, e não o assunto da pergunta. Os autores relatam que os estudantes muitas vezes perguntavam sobre assuntos diferentes do tópico escolhido no aplicativo, o que exigiu a categorização manual das perguntas \parencite{zaman2024dataset}. Ainda assim, 44 perguntas rotuladas como Revisão de \textit{Python}, possivelmente, tratam, pelo texto, de outros tópicos: 25 de recursão, como ``Como resolver o problema da Torre de Hanói usando recursão?”, 6 de grafos, 4 de \textit{hashing} e 9 de outros tópicos. A plataforma atribuiu 43 delas ao tópico do assunto, assim como uma pergunta rotulada como Grafo que trata de árvore binária AVL, tipo de árvore binária de busca balanceada cujo nome vem das iniciais dos sobrenomes de seus criadores, Adelson-Velsky e Landis.
}

\novoBlocoTexto{
    A terceira são termos com mais de um sentido. Em Pilhas, 37 das 196 perguntas tratam da pilha como região de memória (\textit{heap}, memória, estouro de pilha) e 19 citam recursão ou a pilha de chamadas. No material, a memória aparece no tópico Ponteiros e a pilha de chamadas no de Recursão, para os quais foram levadas 9 e 10 perguntas de Pilhas; as 9 levadas a Ponteiros tratam todas da pilha como memória do programa, como a alocação de memória na pilha, o estouro de \textit{buffer} e o registrador de ponteiro de pilha. Do mesmo modo, 7 das 14 perguntas de Nós citam árvore, raiz, folha ou filho, e as 6 perguntas de Nós atribuídas a Árvore binária de busca tratam dos nós de uma árvore.
}

\novoBlocoTexto{
    O Quadro \ref{quad:divergencias-classificacao} apresenta um exemplo de pergunta para cada uma dessas causas.
}

\begin{quadro}[H]\captionsetup{name=Quadro}
\centering
\caption{Exemplos de divergências na classificação das perguntas}
\label{quad:divergencias-classificacao}

\renewcommand{\arraystretch}{1.3}
\setlength{\tabcolsep}{3pt}

\begin{tabularx}{\textwidth}{
|>{\raggedright\arraybackslash}p{2.5cm}
|>{\raggedright\arraybackslash}X
|>{\raggedright\arraybackslash}p{2.2cm}
|>{\raggedright\arraybackslash}p{2.0cm}|
}
\hline

\rowcolor[HTML]{C0C0C0}
\textbf{Causa provável} &
\textbf{Pergunta enviada à plataforma} &
\textbf{Tópico de referência} &
\textbf{Tópico atribuído} \\ \hline

Divisão diferente dos tópicos &
Qual é a diferença entre uma lista e uma tupla em Python, tanto em termos de sintaxe quanto de mutabilidade? &
Revisão de Python &
Listas \\ \hline

Rótulo pelo tema da aula &
Como calcular o Máximo Divisor Comum (MDC) de dois números usando recursão? &
Revisão de Python &
Recursão \\ \hline

Termo com mais de um sentido &
Quais são os benefícios e os desafios da alocação de memória na pilha (\textit{stack}) em comparação com a alocação no \textit{heap}? &
Pilhas &
Ponteiros \\ \hline

\end{tabularx}

{\fontsize{10pt}{12pt}\selectfont
\caption*{Fonte: Elaborado pela autora (2026).}
}
\end{quadro}

\FloatBarrier


\FloatBarrier

\novoBlocoTexto{
A etapa em que cada divergência surgiu, na recuperação dos fragmentos ou na escolha da LLM, não foi quantificada nesta avaliação, embora o banco de dados registre os cinco fragmentos recuperados para cada pergunta. O painel mostra, porém, que a busca também recupera fragmentos de outros tópicos: entre os slides mais recuperados nas dúvidas atribuídas a Pilhas, está um slide do material de Filas (Figura \ref{fig:tela-topico-trechos}).
}

\novoBlocoTexto{
Quanto à tradução, a comparação entre o texto original e o traduzido mostrou que os principais termos técnicos foram traduzidos como no material da disciplina em quase todas as perguntas; linked list, por exemplo, foi traduzido como lista encadeada em 148 das 151 perguntas em que aparece. Traduções incomuns, como “lista duplamente articulada” e “percussão” no lugar de percurso (traversal), ocorreram em 23 perguntas, e só 3 delas foram atribuídas a outro tópico. Em duas, a divergência já se explica pelas causas anteriores, e apenas em uma, sobre a detecção de ciclos em uma “lista duplamente articulada”, atribuída a Grafo, a tradução pode ter contribuído. Assim, nos termos conferidos, a tradução não explica as divergências observadas. A ambiguidade do termo pilha, que designa tanto a estrutura de dados quanto a região de memória, já existe no inglês stack e, portanto, não decorre da tradução.
}

\novoBlocoTexto{
Para o professor, essas divergências aparecem na contagem de dúvidas por tópico do painel (Figura \ref{fig:tela-painel}). Os seis tópicos com mais dúvidas são os mesmos na referência e na plataforma (Revisão de Python, Árvore binária de busca, Pilhas, Recursão, \textit{Hashing} (Espalhamento) e Grafo), mas em outra ordem: Revisão de Python, o tópico com mais perguntas na referência (295), aparece no painel em 6º lugar (137), enquanto Recursão passa de 163 para 206 dúvidas e Listas, de 21 para 85. Assim, o painel indica os conteúdos com mais dúvidas, mas subestima Revisão de Python e superestima Listas e Recursão, o que reforça que ele deve ser lido como subsídio, e não como diagnóstico definitivo (seção~\ref{sec:classificacao-textual-llm}).
}
\subsection{{\fontsize{12pt}{14pt}\selectfont Discussão e limitações}}
\label{sec:discussao-limitacoes}
\novoBlocoTexto{
Os resultados indicam que a plataforma organizou as dúvidas pelos tópicos do material do professor com 83,33\% de concordância com a classificação manual de referência. Em relação ao problema desta pesquisa, isso significa que o professor pode reconhecer quais tópicos concentram mais dúvidas: os seis tópicos com mais perguntas foram os mesmos na referência e na plataforma (seção~\ref{sec:analise-erros}). Essa organização dispensa a categorização manual, tarefa trabalhosa que, na prática, impede o professor sem tempo disponível de aproveitar esses dados \parencite{lobo2025democratizing}.
}

\novoBlocoTexto{
O desempenho foi maior nos tópicos de conteúdo bem delimitado, como \textit{Hashing} (Espalhamento), Árvore binária de busca e Grafo, com F1 acima de 92\%, e menor nos tópicos cujo conteúdo se sobrepõe ao de outros, como Revisão de Python, Listas, Nós e Ponteiros. Parte dessas divergências decorre do critério de classificação: a referência agrupa parte das perguntas pelo tema da aula em que foram feitas, enquanto a plataforma as classifica pelo conteúdo do material do professor, que é a proposta deste trabalho. Assim, uma pergunta de Revisão de Python sobre listas atribuída ao tópico Listas conta como erro na avaliação, mas aponta ao professor a parte do material relacionada à dúvida, como ocorreu com 39 das 63 perguntas de Revisão de Python atribuídas a Listas e com 44 das 45 perguntas cujo rótulo difere do assunto (seção~\ref{sec:analise-erros}).
}

\novoBlocoTexto{
Esse resultado deve ser lido considerando a dificuldade da própria tarefa. No estudo de \textcite{qiu2025systematic}, em uma avaliação humana de um subconjunto de 285 itens rotulados por uma LLM, três avaliadores independentes concordaram totalmente em 83,16\% dos itens (237 de 285), com kappa de 0,462, valor que os autores classificam como concordância moderada. Isso sugere que, mesmo entre avaliadores humanos, o julgamento do tópico de uma dúvida não é unânime. \textcite{lobo2025democratizing}, por sua vez, observaram desempenho fraco de modelos com cerca de 3 e 3,8 bilhões de parâmetros e lembram que, no estudo de \textcite{goldberg2021refining}, a concordância entre professores ao rotular perguntas foi de apenas cerca de 65\%; neste trabalho, um modelo de 7 bilhões de parâmetros, quantizado e executado localmente, com apoio do RAG, chegou a 83,33\%. Essas comparações são apenas indicativas, pois as tarefas diferem: em \textcite{qiu2025systematic} e em \textcite{goldberg2021refining}, os valores medem a concordância entre avaliadores, e não entre a LLM e uma referência, e em \textcite{lobo2025democratizing}, classificava-se o tipo da pergunta, e não o tópico.
}

\novoBlocoTexto{
Mesmo com esse nível de concordância, uma em cada seis perguntas (16,67\%) foi apresentada em um tópico diferente do de referência, e as divergências, concentradas em poucos pares de tópicos, alteraram a ordem dos tópicos no painel (seção~\ref{sec:analise-erros}). Por isso, como discutido na seção~\ref{sec:classificacao-textual-llm}, a contagem por tópico deve ser apresentada ao professor como subsídio à análise pedagógica, e não como diagnóstico definitivo, cabendo a ele interpretar os dados e definir as ações, etapas que \textcite{wise2019teaching} apontam como parte da tomada de decisão docente fundamentada em dados.
}

\novoBlocoTexto{
O experimento tem limitações. As perguntas foram escritas em inglês e traduzidas automaticamente, parte dos rótulos de referência indica o tema da aula, e não o assunto da pergunta; os próprios autores reconhecem que o conjunto contém perguntas vagas ou mal estruturadas \parencite{zaman2024dataset}. Os itens do conteúdo programático foram gerados a partir dos tópicos extraídos do conjunto de referência, e as descrições dos tópicos extraídas do plano são enviadas à LLM na classificação, o que pode ter favorecido a concordância obtida (seção~\ref{sec:preparacao-disciplina}). Além disso, avaliou-se uma única disciplina, em uma única execução, com um único modelo e uma única configuração, sem medir a variação entre execuções, e as perguntas foram enviadas por automação, a partir de uma única conta, e não por estudantes de uma turma real. A etapa em que surgiu cada divergência não foi quantificada. Por fim, o material da disciplina foi elaborado a partir de apostilas para o experimento, e não pelo professor da disciplina, e a qualidade das respostas e dos resumos não foi avaliada.
}
